% Generated draft: requires mathematical and source review.
\documentclass[11pt]{article}
\usepackage[margin=25mm]{geometry}
\usepackage{fontspec}
\setmainfont{Latin Modern Roman}
\usepackage{amsmath,amssymb,mathtools,unicode-math}
\setmathfont{Latin Modern Math}
\usepackage{enumitem,needspace,xurl,hyperref}
\setlength{\parindent}{0pt}
\setlength{\parskip}{5pt}
\setlength{\emergencystretch}{3em}
\newcommand{\N}{\mathbb{N}}
\newcommand{\Z}{\mathbb{Z}}
\newcommand{\Q}{\mathbb{Q}}
\newcommand{\R}{\mathbb{R}}
\newcommand{\C}{\mathbb{C}}
\newcommand{\F}{\mathbb{F}}
\newcommand{\A}{\mathbb{A}}
\newcommand{\PP}{\mathbb{P}}
\newcommand{\E}{\mathbb{E}}
\newcommand{\eps}{\varepsilon}
\newcommand{\dd}{\,\mathrm{d}}
\begin{document}
\section*{UnsolvedMath ID: 7600020}
\textbf{Source problem code:} catalogue entry
\subsection{Definitions and mathematical statement}
Monadic second-order logic of an order has variables for individual points and for unary subsets, with order and membership relations. The Borel monadic theory of the real order (BMTO) is the set of sentences true in $(\mathbb R,\leq)$ when every set quantifier ranges over Borel subsets of $\mathbb R$. Is this set of sentence codes decidable, meaning that a single algorithm halts on every sentence and correctly decides its truth?

The second question asks whether the monadic theory of well-ordering (MTWO), with unrestricted subset quantifiers, is consistently decidable. A usual class-theory reading takes
\[
\operatorname{MTWO}=\{\varphi:(\alpha,<)\models_{\mathrm{MSO}}\varphi\text{ for every ordinal }\alpha\},
\]
with full power-set semantics for each ordinal. Under that reading, the consistency question asks whether the intended background set theory can consistently hold that this theory is recursive. The source does not specify the background theory or distinguish the theory of all well-orders from that of particular well-orders; this interpretation must be checked rather than silently identifying the variants.

The supplied summary reports BMTO decidable in 2024--2025 work, while keeping MTWO's consistency-decidability question open. The two set-quantifier ranges and the direct versus consistency questions are separate.
\subsection{Short English statement}
Decide the Borel-set monadic theory of the real order, and separately ask whether the unrestricted monadic theory of well-ordering can consistently be decidable. The supplied summary reports the first part solved and leaves the second open.
\subsection{Sources}
\begin{itemize}
\item Wikipedia contributors, \emph{List of unsolved problems in mathematics}, ``Model theory and formal languages,'' bullet 20, pinned revision 1366636767, \url{https://en.wikipedia.org/w/index.php?title=List_of_unsolved_problems_in_mathematics&oldid=1366636767#Model_theory_and_formal_languages}; accessed 2026-07-29, MediaWiki source extraction. Wikipedia text is CC BY-SA 4.0, with attribution/share-alike handling required for release.
\item Manthe, ``The Borel monadic theory of order is decidable,'' arXiv:2410.00887. The supplied summary attributes this to a resolution of Shelah's conjecture and reports an elementary substructure given by Boolean combinations of $F_\sigma$ sets.
\item ``MSO logic of the real order with set quantifiers ranging over the Borel sets,'' arXiv:2512.23003; authors are not supplied. The summary reports weak/strong forms and references interpretations in S2S and B\"uchi-style decidability machinery, without their precise hypotheses.
\item Local registry \texttt{data/unsolvedmath/problems.json}, permanent UnsolvedMath ID $7600020$, source code catalogue entry, same two-part titled problem. Dataset attribution: Ulam AI and the UnsolvedMath Contributors.
\end{itemize}
\end{document}
